%=============================================================================
\section{Programování v Javě a objektově-orientované paradigma}
%=============================================================================

\subsection{Objektově-orientované programování (OOP)}

\subsubsection{Základní principy OOP}

Objektově-orientované programování je programovací paradigma organizující software kolem objektů,
které kombinují data (atributy/pole) a chování (metody).

\textbf{Lidsky řečeno:} Místo toho, abys psal program jako jeden dlouhý seznam příkazů, OOP tě nutí přemýšlet o světě jako o věcech (objektech), které mají vlastnosti a umějí něco dělat. Auto má značku, barvu a počet dveří (data) -- a umí nastartovat, zastavit a zatočit (metody). V programu pak vytvoříš třídu \texttt{Auto} jako šablonu a z ní libovolně mnoho konkrétních aut (objektů). Každé auto si pamatuje své vlastní hodnoty, ale sdílí stejné chování. Díky tomu je kód přehledný, znovupoužitelný a snáz se v něm orientuje víc lidí.

\paragraph{1. Abstrakce (Abstraction)}
Proces skrytí složitých implementačních detailů a zobrazení pouze nezbytné funkcionality.
\begin{itemize}
  \item Zaměřuje se na \textit{co} objekt dělá, ne \textit{jak} to dělá
  \item Realizuje se pomocí abstraktních tříd (\texttt{abstract class}) a rozhraní (\texttt{interface})
  \item Příklad: rozhraní \texttt{List} v Javě -- programátor pracuje s metodami \texttt{add()}, \texttt{get()},
    \texttt{remove()}, aniž ví, zda jde o \texttt{ArrayList} nebo \texttt{LinkedList}
\end{itemize}

\begin{lstlisting}[language=Java, caption=Abstrakce -- rozhraní Tvar]
// Abstrakce: definujeme CO umí tvar, ne JAK to dělá
interface Tvar {
    double obsah();       // kazdý tvar musí umět spočítat obsah
    double obvod();
    String nazev();
}

// Konkrétní implementace -- SKRYTÉ od uživatele
class Kruh implements Tvar {
    private double polomer;
    public Kruh(double r) { this.polomer = r; }

    @Override public double obsah() { return Math.PI * polomer * polomer; }
    @Override public double obvod() { return 2 * Math.PI * polomer; }
    @Override public String nazev()  { return "Kruh"; }
}

class Obdelnik implements Tvar {
    private double a, b;
    public Obdelnik(double a, double b) { this.a = a; this.b = b; }

    @Override public double obsah() { return a * b; }
    @Override public double obvod() { return 2 * (a + b); }
    @Override public String nazev()  { return "Obdelnik"; }
}

// Klientský kód pracuje POUZE s abstrakcí Tvar
// nezajímá ho, co je uvnitř
public static void vypisTvary(List<Tvar> tvary) {
    for (Tvar t : tvary) {
        System.out.printf("%s: obsah=%.2f%n", t.nazev(), t.obsah());
    }
}
\end{lstlisting}

\textbf{Proč je to abstrakce?} Funkce \texttt{vypisTvary} neví, zda pracuje s kruhem nebo obdélníkem -- pracuje pouze s rozhraním \texttt{Tvar}. Implementační detaily (vzorec pro obsah kruhu) jsou skryté za rozhraním. Lze přidat nový tvar (\texttt{Trojuhelnik}) \textbf{bez jakékoliv změny} v \texttt{vypisTvary}.

\paragraph{2. Zapouzdření (Encapsulation)}
Sdružení dat a metod, které s nimi pracují, do jednoho celku (třídy) a omezení přístupu zvenčí.
\begin{itemize}
  \item Data (atributy) by měla být \texttt{private}, přístup přes \texttt{public} gettery/settery
  \item Chrání interní stav objektu před neplatným nastavením
  \item \textbf{Modifikátory přístupu v Javě:}
    \begin{itemize}
      \item \texttt{private} -- přístup pouze z dané třídy
      \item \texttt{protected} -- přístup z třídy a podtříd (\textbf{podtřída} = třída, která dědí z jiné třídy; v Javě \texttt{class Pes extends Zvire}, v C\# přesně \texttt{class Pes : Zvire} -- syntaxe se liší, princip je stejný)
      \item \texttt{package-private} (výchozí) -- přístup v rámci balíčku
      \item \texttt{public} -- přístup odkudkoliv
    \end{itemize}
\end{itemize}

\begin{lstlisting}[language=Java, caption=Zapouzdření v Javě]
public class BankovniUcet {
    private double zustatek;  // soukromé pole

    public double getZustatek() {
        return zustatek;
    }

    public void vlozit(double castka) {
        if (castka > 0) zustatek += castka;
        // validace vstupu pred zmenou stavu
    }
}
\end{lstlisting}

\paragraph{3. Dědičnost (Inheritance)}
Mechanismus, který umožňuje nové třídě přebírat vlastnosti a metody existující (nadřazené) třídy.
\begin{itemize}
  \item Klíčové slovo \texttt{extends} (třída) nebo \texttt{implements} (rozhraní) --
    rozdíl: \texttt{extends} použiješ, když dědíš z \textbf{konkrétní nebo abstraktní třídy}
    (přebíráš hotový kód, atributy, chování); \texttt{implements} použiješ, když třída
    plní \textbf{rozhraní} (interface) -- rozhraní je jen „smlouva" co třída musí umět,
    bez jakékoliv implementace. V C\# se obojí píše jako \texttt{:}, Java to rozlišuje
    syntakticky.
  \item Java podporuje jednoduchou dědičnost tříd (jen jeden \texttt{extends}), ale vícenásobnou implementaci rozhraní (\texttt{implements A, B, C})
  \item \texttt{super} -- odkaz na nadřazenou třídu
  \item \textbf{IS-A} vztah (pes IS-A zvíře) -- říká, že podtřída \textit{je druhem} nadtřídy;
    dědičnost je správná právě tehdy, když platí IS-A: \texttt{Pes extends Zvire} dává smysl,
    protože pes opravdu \textit{je} zvíře. Pokud IS-A neplatí, dědičnost nepoužívej.
  \item \textbf{HAS-A} vztah (auto HAS-A motor) -- kompozice; třída \textit{obsahuje} jinou třídu
    jako atribut, místo aby z ní dědila. \texttt{Auto} nemá \textit{být} motor, ale \textit{mít} ho
    (\texttt{private Motor motor;}). Kompozice je obecně preferována před dědičností -- méně těsné propojení.
  \item \textbf{Nevýhody dědičnosti:} těsné propojení (coupling) tříd; změna v rodičovské třídě může rozbít potomky
\end{itemize}

\paragraph{4. Polymorfismus (Polymorphism)}
Schopnost různých objektů reagovat různým způsobem na stejné zprávy (volání metod).
\begin{itemize}
  \item \textbf{Compile-time polymorfismus} -- přetížení metod (method overloading); stejný název metody,
    různé parametry
  \item \textbf{Runtime polymorfismus} -- přepsání metody (method overriding); podtřída překryje metodu
    nadtřídy; rozhodnutí za běhu programu
  \item \textbf{Liskovův substituční princip (LSP):} objekty podtřídy musejí být zaměnitelné za objekty nadtřídy
\end{itemize}

\begin{lstlisting}[language=Java, caption=Polymorfismus v Javě]
abstract class Zvire {
    abstract void vydejZvuk();
}

class Pes extends Zvire {
    @Override
    public void vydejZvuk() { System.out.println("Haf!"); }
}

class Kocka extends Zvire {
    @Override
    public void vydejZvuk() { System.out.println("Mnau!"); }
}

// Runtime polymorfismus
Zvire z = new Pes();
z.vydejZvuk(); // vypise "Haf!" -- rozhodnutí za běhu
\end{lstlisting}

\subsubsection{Třída vs. instance}

\begin{itemize}
  \item \textbf{Třída (Class)} -- šablona/plán; definuje strukturu a chování; na disku jako \texttt{.class} soubor
  \item \textbf{Instance (Object)} -- konkrétní výskyt třídy vytvořený v paměti pomocí \texttt{new}
  \item \textbf{Konstruktor} -- speciální metoda volaná při vytváření instance; stejný název jako třída, bez návratového typu
\end{itemize}

\subsubsection{Abstraktní třídy vs. rozhraní}

\paragraph{Co je abstraktní třída?}
\textbf{Abstraktní třída} je třída, od které \textbf{nelze přímo vytvořit instanci} (\texttt{new} nad ní selže).
Slouží jako \textbf{neúplná šablona} -- definuje společné chování potomků, ale části nechává k implementaci v podtřídách.
Používá se, když třídy sdílejí \textbf{implementační kód} i \textbf{stav} (atributy).

\begin{lstlisting}[language=Java, caption=Abstraktní třída]
abstract class Vozidlo {
    protected String znacka;     // sdílený stav
    protected int rokVyroby;

    public Vozidlo(String znacka, int rok) {
        this.znacka = znacka;
        this.rokVyroby = rok;
    }

    // Konkrétní metoda -- sdílena všemi potomky
    public int stariVozu() {
        return 2024 - rokVyroby;
    }

    // Abstraktní metoda -- každý potomek musí implementovat
    public abstract String typPaliva();
}

class Auto extends Vozidlo {
    public Auto(String znacka, int rok) { super(znacka, rok); }
    @Override public String typPaliva() { return "Benzin/nafta/elektro"; }
}
\end{lstlisting}

\paragraph{Co je rozhraní (interface)?}
\textbf{Rozhraní} je \textbf{čistý kontrakt} -- říká ``tato třída UMĚT dělat X'', ale vůbec neříká jak.
Neobsahuje stav (žádné instanční proměnné). Třída může implementovat \textbf{libovolný počet} rozhraní.
Používá se pro definování schopností nezávislých na hierarchii tříd.

\begin{lstlisting}[language=Java, caption=Rozhraní]
interface Elektricke {
    void nabijBaterii();
    int getKapacitaBaterie();  // kWh
}

interface GPS {
    double getSirka();
    double getDelka();
}

// Třída může implementovat více rozhraní zároveň
class ElektrickeAuto extends Auto implements Elektricke, GPS {
    private int kapacita;
    private double lat, lon;

    @Override public void nabijBaterii() { /* ... */ }
    @Override public int getKapacitaBaterie() { return kapacita; }
    @Override public double getSirka() { return lat; }
    @Override public double getDelka() { return lon; }
}
\end{lstlisting}

\paragraph{Kdy použít co?}

\begin{center}
\begin{tabular}{p{5.5cm}p{5.5cm}}
\toprule
\textbf{Abstraktní třída} & \textbf{Rozhraní} \\
\midrule
Třídy \textbf{jsou} stejným typem (IS-A): \newline
\texttt{Pes} IS-A \texttt{Zvire} &
Třídy \textbf{umí} totéž (CAN-DO): \newline
\texttt{Auto} CAN-DO \texttt{Serializable} \\
\midrule
Sdílí \textbf{kód} i \textbf{stav} (atributy) &
Definují pouze \textbf{kontrakt} (co umí) \\
\midrule
Potřebujeme \textbf{konstruktor} nebo \textbf{protected} pole &
Přidáváme schopnosti \textbf{napříč hierarchiemi} \\
\midrule
Pouze \textbf{jedna} abstraktní třída &
\textbf{Více} rozhraní najednou \\
\bottomrule
\end{tabular}
\end{center}

\subsubsection{Principy SOLID}

\begin{itemize}
  \item \textbf{S} -- Single Responsibility Principle: třída má mít pouze jeden důvod ke změně.
    Každá třída dělá jednu věc a dělá ji dobře. Třída \texttt{Invoice} počítá cenu faktury --
    ale tisk faktury do PDF má mít vlastní třídu \texttt{InvoicePrinter}. Pokud by se změnil
    formát PDF, nemusíš sahat do logiky výpočtu ceny.

  \item \textbf{O} -- Open/Closed Principle: třída je otevřena pro rozšíření, ale uzavřena pro
    modifikaci. Přidáš nové chování přes nové třídy (dědičností nebo kompozicí), ne tím, že
    přepisuješ stávající kód. Příklad: místo přidání \texttt{if (typ == "kruh")} do existující
    metody vytvoříš novou třídu \texttt{Kruh}, která implementuje rozhraní \texttt{Tvar}.

  \item \textbf{L} -- Liskov Substitution Principle: podtřídy musejí být plně zaměnitelné za
    nadtřídy -- kód, který pracuje s \texttt{Zvire}, musí fungovat stejně správně i když mu
    předáš \texttt{Pes} nebo \texttt{Kocka}. Pokud podtřída mění nebo ruší chování rodičovské
    třídy (háže výjimku tam, kde rodič nic nedělá), porušuje LSP.

  \item \textbf{I} -- Interface Segregation Principle: raději víc malých specifických rozhraní
    než jedno velké „all-in-one". Třída by neměla být nucena implementovat metody, které
    nepotřebuje. Místo jednoho rozhraní \texttt{Stroj} s metodami \texttt{tiskni()},
    \texttt{skenuj()}, \texttt{faxuj()} udělej tři samostatná -- tiskárna pak implementuje
    jen \texttt{Tiskatelny}, ne zbylé dvě.

  \item \textbf{D} -- Dependency Inversion Principle: třídy vysoké úrovně by neměly záviset
    na třídách nízké úrovně -- oboje by mělo záviset na abstrakcích (rozhraních). Příklad:
    \texttt{ObjednavkoveService} nemá přímo vytvářet \texttt{new MySQLDatabaze()}, ale dostat
    rozhraní \texttt{Databaze} zvenčí (injekce závislostí). Snadno tak vyměníš MySQL za
    jinou DB bez dotyku servisní logiky.
\end{itemize}

%=============================================================================
\subsection{Test Driven Development (TDD)}
%=============================================================================

\subsubsection{Principy TDD}

\textbf{TDD} je vývojová metodika, při níž se testy píší \textbf{před} samotným produkčním kódem.
Cyklus TDD: \textbf{Red -- Green -- Refactor}

\begin{enumerate}
  \item \textbf{Red} -- napsat test, který selže (testuje chování, které ještě neexistuje)
  \item \textbf{Green} -- napsat minimální kód tak, aby test prošel
  \item \textbf{Refactor} -- vylepšit kód bez změny chování (testy stále musí procházet)
\end{enumerate}

\subsubsection{Typy testů}

\begin{itemize}
  \item \textbf{Unit testy} -- testují izolovanou jednotku (třídu, metodu); mock objekty pro závislosti
  \item \textbf{Integrační testy} -- testují spolupráci více komponent (např. DB + služba)
  \item \textbf{End-to-end (E2E) testy} -- testují celý systém z pohledu uživatele
\end{itemize}

\textbf{Testovací pyramida:} mnoho unit testů, méně integračních, nejméně E2E.

\subsubsection{JUnit 5 v Javě}

\begin{lstlisting}[language=Java, caption=Unit test s JUnit 5]
import org.junit.jupiter.api.*;
import static org.junit.jupiter.api.Assertions.*;

class KalkulackaTest {
    private Kalkulacka calc;

    @BeforeEach
    void setUp() {
        calc = new Kalkulacka();
    }

    @Test
    void testScitani() {
        assertEquals(5, calc.secti(2, 3));
    }

    @Test
    void testDeleniNulou() {
        assertThrows(ArithmeticException.class,
            () -> calc.vydel(10, 0));
    }
}
\end{lstlisting}

\subsubsection{Mockování závislostí}

Mock objekty simulují chování závislostí (DB, externích API) při unit testech.
Knihovna \textbf{Mockito} v Javě:
\begin{lstlisting}[language=Java, caption=Mockování s Mockito]
@Mock
private UzivatelskyRepository repo;

@Test
void testNajdiUzivatele() {
    when(repo.findById(1L))
        .thenReturn(Optional.of(new Uzivatel("Jan")));
    Uzivatel u = service.najdi(1L);
    assertEquals("Jan", u.getJmeno());
    verify(repo).findById(1L);
}
\end{lstlisting}

%=============================================================================
\subsection{Lambda výrazy a Stream API}
%=============================================================================

\subsubsection{Funkcionální rozhraní a lambda výrazy}

\textbf{Lambda výraz} je anonymní funkce -- blok kódu, který lze předat jako argument nebo uložit do proměnné.
Zavedeny v Javě 8. Syntaxe: \texttt{(parametry) -> \{tělo\}}

\textbf{Lidsky:} Před Javou 8 jsi musel pro každou „malou akci" psát celou anonymní třídu (viz příklad níže).
Lambda je zkratka -- místo tří řádků boilerplate napíšeš jednu řádku. Jde vlastně o způsob, jak předat
\textit{chování} (funkci) jako hodnotu, stejně jako předáváš číslo nebo řetězec.

\textbf{Funkcionální rozhraní} -- rozhraní s právě jednou abstraktní metodou (označuje se \texttt{@FunctionalInterface}).
Lambda výraz \textit{je} implementace takového rozhraní -- Java ví, kterou jedinou metodu má lambda naplnit.

\textbf{Zabudovaná funkcionální rozhraní} -- Java 8 dodává sadu hotových rozhraní pro nejčastější případy,
abys nemusel psát vlastní. Tabulka níže ukazuje co každé „očekává" a co „vrací":

\begin{center}
\begin{tabular}{llll}
\toprule
\textbf{Rozhraní} & \textbf{Vstup $\to$ Výstup} & \textbf{Použití} & \textbf{Příklad} \\
\midrule
\texttt{Function<T,R>} & T $\to$ R & Transformace hodnoty & \texttt{s -> s.length()} \\
\texttt{Predicate<T>} & T $\to$ boolean & Filtrování (ano/ne) & \texttt{s -> s.isEmpty()} \\
\texttt{Consumer<T>} & T $\to$ nic & Spotřebování, vedlejší efekt & \texttt{s -> System.out.println(s)} \\
\texttt{Supplier<T>} & nic $\to$ T & Produkce hodnoty & \texttt{() -> new ArrayList<>()} \\
\texttt{BiFunction<T,U,R>} & T, U $\to$ R & Transformace dvou hodnot & \texttt{(a,b) -> a + b} \\
\bottomrule
\end{tabular}
\end{center}

\textbf{Jak číst tabulku:} \texttt{Function} vezme jednu věc a vrátí jinou (např. ze \texttt{String} udělá \texttt{Integer} -- délku).
\texttt{Predicate} vezme věc a řekne pravda/nepravda -- proto se používá ve \texttt{filter()}.
\texttt{Consumer} věc vezme, ale nic nevrátí -- jen ji ``spotřebuje'' (vypíše, uloží...).
\texttt{Supplier} naopak nic nebere, ale hodnotu \textit{vyrobí}.
\texttt{BiFunction} je jako \texttt{Function}, ale bere dva vstupy místo jednoho.

\begin{lstlisting}[language=Java, caption=Lambda výrazy]
// Starý způsob - anonymní třída
Comparator<String> comp1 = new Comparator<String>() {
    public int compare(String a, String b) {
        return a.compareTo(b);
    }
};

// Lambda výraz
Comparator<String> comp2 = (a, b) -> a.compareTo(b);

// Method reference
Comparator<String> comp3 = String::compareTo;
\end{lstlisting}

\subsubsection{Stream API}

\textbf{Stream} je \textit{pipeline} pro zpracování kolekce dat funkcionálním stylem -- místo imperativního cyklu \texttt{for} popisujeme, co chceme s daty udělat jako řetěz operací.

\textbf{Klíčové vlastnosti:}
\begin{itemize}
  \item \textbf{Lazy evaluation} -- intermediární operace (\texttt{filter}, \texttt{map}~\ldots) se
    \textit{nevyhodnocují hned} -- jen popíšou, co se má stát. Skutečné zpracování proběhne až
    ve chvíli, kdy přijde terminální operace (\texttt{collect}, \texttt{findFirst}~\ldots).
    \textbf{Proč je to dobré?} Představ si seznam milionu čísel a chceš první sudé větší než 100.
    Bez lazy evaluation bys profiltroval všech milion prvků a pak vzal první výsledek.
    S lazy evaluation stream okamžitě zastaví, jakmile najde první shodu -- zbytek vůbec nezpracuje.

  \item \textbf{Spotřeba} -- stream lze projít pouze jednou; po terminální operaci je ``vyčerpaný''

  \item \textbf{Nemodifikuje zdroj} -- stream nikdy nemění původní kolekci, ze které vznikl;
    \texttt{filter()} nebo \texttt{map()} vždy vrátí \textit{nový} stream s novými daty.
    Původní \texttt{List} zůstane nedotčený.
    \textbf{Příklad:} zavoláš \texttt{seznam.stream().filter(...).collect(...)} --
    \texttt{seznam} má pořád stejné prvky jako předtím, výsledek filtrování je v nové proměnné.
  \item \textbf{Paralelizovatelný} -- \texttt{parallelStream()} automaticky rozdělí práci mezi vlákna (Fork/Join pool)
\end{itemize}

\paragraph{Operace na streamech}
\begin{itemize}
  \item \textbf{Intermediární (lazy)} -- vracejí nový stream: \texttt{filter()}, \texttt{map()}, \texttt{flatMap()},
    \texttt{sorted()}, \texttt{distinct()}, \texttt{limit()}, \texttt{peek()}

    \textbf{Rozdíl \texttt{map} vs \texttt{flatMap}:}
    \texttt{map()} transformuje každý prvek na jeden výsledek -- ze streamu N prvků vznikne stream N prvků.
    \texttt{flatMap()} transformuje každý prvek na \textit{stream} výsledků a pak všechny tyto streamy
    „zploští" do jednoho -- použiješ ho, když by \texttt{map()} vrátil stream streamů.
    Příklad: máš seznam vět a chceš stream jednotlivých slov --
    \texttt{map} by vrátil \texttt{Stream<String[]>} (každá věta jako pole),
    \texttt{flatMap} vrátí rovnou \texttt{Stream<String>} (všechna slova v jednom streamu):
\begin{lstlisting}[language=Java]
List<String> vety = List.of("ahoj svete", "jak se mas");

// map -- Stream<String[]>, vnořené pole
vety.stream().map(v -> v.split(" "));

// flatMap -- Stream<String>, plochý seznam slov
vety.stream()
    .flatMap(v -> Arrays.stream(v.split(" ")))
    .collect(Collectors.toList());
// ["ahoj", "svete", "jak", "se", "mas"]
\end{lstlisting}
  \item \textbf{Terminální (eager)} -- vrátí výsledek nebo vedlejší efekt: \texttt{collect()}, \texttt{forEach()},
    \texttt{reduce()}, \texttt{count()}, \texttt{findFirst()}, \texttt{anyMatch()}, \texttt{allMatch()}
\end{itemize}

\begin{lstlisting}[language=Java, caption=Stream API příklady]
List<String> jmena = Arrays.asList("Alice", "Bob", "Charlie", "David");

// Filtrování a transformace
List<String> vysledek = jmena.stream()
    .filter(j -> j.length() > 3)      // Alice, Charlie, David
    .map(String::toUpperCase)          // ALICE, CHARLIE, DAVID
    .sorted()                          // řazení
    .collect(Collectors.toList());

// Redukce
int soucet = IntStream.rangeClosed(1, 10)
    .reduce(0, Integer::sum);   // 55

// Seskupení
Map<Integer, List<String>> podleDlky = jmena.stream()
    .collect(Collectors.groupingBy(String::length));

// Paralelní stream (pro velká data)
long pocet = jmena.parallelStream()
    .filter(j -> j.startsWith("A"))
    .count();
\end{lstlisting}

\paragraph{Srovnání Java Stream API s C\# LINQ}

Pokud znáte .NET, Stream API je ekvivalentem \textbf{LINQ} (Language Integrated Query).
Koncepty jsou totožné, liší se pouze syntaxí:

\begin{lstlisting}[language=Java, caption=Java Stream API vs. C\# LINQ -- srovnání]
// === JAVA (Stream API) ===
List<String> jmena = List.of("Alice", "Bob", "Charlie", "David");

List<String> vysledek = jmena.stream()
    .filter(j -> j.length() > 3)    // jako Where() v C#
    .map(String::toUpperCase)        // jako Select() v C#
    .sorted()
    .collect(Collectors.toList());   // jako ToList() v C#

// Průměrný věk dospělých
double prumer = osoby.stream()
    .filter(o -> o.getVek() >= 18)
    .mapToInt(Osoba::getVek)         // speciální IntStream pro int
    .average()
    .orElse(0.0);

// === C# LINQ (ekvivalent) ===
// var vysledek = jmena
//     .Where(j => j.Length > 3)
//     .Select(j => j.ToUpper())
//     .OrderBy(j => j)
//     .ToList();
//
// double prumer = osoby
//     .Where(o => o.Vek >= 18)
//     .Average(o => o.Vek);
\end{lstlisting}

\textbf{Hlavní rozdíly Java Stream vs. C\# LINQ:}
\begin{itemize}
  \item V Javě je \texttt{stream()} explicitní volání; v C\# je LINQ přímo na \texttt{IEnumerable<T>}
  \item Java rozlišuje \texttt{Stream<T>}, \texttt{IntStream}, \texttt{LongStream}, \texttt{DoubleStream} (kvůli boxing); C\# to řeší automaticky
  \item Terminální operace Javy: \texttt{collect(Collectors.toList())}; v C\# jednoduše \texttt{.ToList()}
  \item Obě platformy podporují paralelní zpracování: Java \texttt{parallelStream()}, C\# \texttt{AsParallel()} (PLINQ)
\end{itemize}

\subsubsection{Optional}

\texttt{Optional<T>} je kontejner pro hodnotu, která může být null -- eliminuje \texttt{NullPointerException}.

\begin{lstlisting}[language=Java, caption=Optional v Javě]
Optional<String> opt = Optional.ofNullable(getJmeno());

// Bezpečná práce s hodnotou
String jmeno = opt.orElse("Neznámý");
opt.ifPresent(j -> System.out.println("Ahoj, " + j));
String velke = opt.map(String::toUpperCase).orElse("");
\end{lstlisting}

\begin{profesorbox}[Otázky profesorů k tématu: Programování v Javě]
\textbf{Komise: Svátek, Chlapek, Zbořil}
\begin{itemize}
  \item Vysvětlete principy OOP -- abstrakce, zapouzdření, dědičnost, polymorfismus.
  \item Jaký je rozdíl mezi abstraktní třídou a rozhraním?
\end{itemize}
\textbf{Vojíř Stanislav, Ing., Ph.D.}
\begin{itemize}
  \item Co je TDD a jaký je jeho postup?
  \item Jak byste testoval metodu, která závisí na externím API?
\end{itemize}
\textbf{Chudán David, Ing., Ph.D.}
\begin{itemize}
  \item Co jsou lambda výrazy a jaké přinášejí výhody?
  \item Vysvětlete rozdíl mezi \texttt{map()} a \texttt{flatMap()} v Stream API.
\end{itemize}
\textbf{Sládek Pavel, Ing., Ph.D.}
\begin{itemize}
  \item Jaké jsou principy SOLID a proč jsou důležité?
  \item Co je polymorfismus a jak se projevuje v Javě?
\end{itemize}
\textbf{Zeman Václav, Ing., Ph.D.}
\begin{itemize}
  \item Jak se v Javě zajišťuje zapouzdření? Jaká jsou pravidla pro modifikátory přístupu?
  \item Proč je dědičnost považována za nebezpečnou a jaká je alternativa?
  \item Co je návrhový vzor publish-subscribe (observer) a proč se využívá?
  \item Co je loose coupling a proč je žádoucí?
  \item Co je Apache Kafka a k čemu slouží?
  \item Co jsou microservices a jaké jsou jejich výhody oproti monolitu?
\end{itemize}
\textbf{Pecinovský Rudolf, RNDr., Ph.D.}
\begin{itemize}
  \item Co je dědičnost a jak se projevuje v OOP?
  \item Jaký je rozdíl mezi nominálním a strukturálním děděním?
\end{itemize}
\textbf{Smutný Zdeněk, Ing., Ph.D.}
\begin{itemize}
  \item Vysvětlete FOR cyklus -- jak funguje a kdy ho použít?
\end{itemize}
\end{profesorbox}
